\documentclass[a4paper]{article}
\usepackage[pdftex]{graphicx}
\usepackage[utf8]{inputenc}
\usepackage{enumerate}
\usepackage{siunitx}
\sisetup{locale=DE}
\DeclareSIUnit \parsec {pc}
\usepackage{href-ul}
\hypersetup{
	colorlinks=true,
	linkcolor=blue,
	urlcolor=blue}
\usepackage{icomma}
\usepackage{amssymb}
\usepackage{geometry}
\geometry{a4paper, top=15mm, left=15mm, right=15mm, bottom=15mm,
	headsep=10mm, footskip=12mm}
\hyphenation{orbit-radius}


\begin{document}
	%Title
	{\bf School assignment from physics }\\
	\begin{enumerate}[1.]
		
		%Task 1
		{\bf \item The planet Neptune has eight known moons, of which only the two outer ones, Triton and Nereid, have been discovered from Earth.} The largest moon Triton, observed in 1846 by the English amateur astronomer William Lassell, requires $T_T= to orbit 5.88$ days with an average orbit radius of $r_T= \SI{354000}{\kilo\meter}$. \\
		Nereid, discovered in 1949 by Dutch astronomer Gerard Kuiper, has an orbital period of $T_N=360.1 $ days
		
		\begin{enumerate}[a)]
			\item Calculate the average orbital speed of the moon Triton in $\SI{}{\kilo\meter\over \second}$
			\item Find the mean orbital radius of Nereid.
		\end{enumerate}
		
		%Exercise 2
		{\bf \item Explain, using Kepler's second law and the technical terms perihelion and aphelion,} why the winter half-year in the northern hemisphere is a few days shorter than the summer half-year.
		
		%Task 3
		{\bf \item Name all astronomical observation results known to you that are considered to be evidence of the correctness of the Big Bang theory.}
		
		%Task 4
		{\bf \item The galaxy NGC 3738}, discovered in 1789 by the German astronomer Wilhelm Herschel, is moving away from us at a speed of $v_G=\SI{331}{\kilo\meter\over \second}$. Calculate the distance of this galaxy in Mpc and in km.\\
		$[\textnormal{Hubble constant:}~H_0=\SI{74}{\kilo\meter\over \second\cdot\mega\parsec}; \quad 1~\textnormal{Parsec~pc}=\SI{3,1e13}{\kilo\meter}]$
		
		%Task 5
		{\bf \item The Triangle Nebula M 33 belongs to the local nebula group of the Milky Way.} It is at a distance of $3.0 \cdot 10^6 $ light-years from it and is moving at a speed of $\SI{190}{\kilo \meter\over \second}$ towards them.
		
		\begin{enumerate}[a)]
			\item In how many years will the Milky Way and M 33 mix?
			\item Give a reason why M 33 is not moving away from us.
		\end{enumerate}
		
		%Task 6
		{\bf \item The diagram opposite shows the schematic side view of our galaxy, the Milky Way.}\\
		
		\begin{minipage}{0.4\textwidth}
			Label them with the following technical terms:\\
			
			discus disk\\
			Halo \\
			Dust layer \\
			Central bulge \\
			Core \\
			Position Sun \\
			Globular clusters.
		\end{minipage}
		\hfill
		\begin{minipage}{0.55\textwidth}
			\includegraphics[width=7 cm]{milchstrasse123.pdf}
		\end{minipage}
		
		
		%Task 7
		{\bf \item Correct any errors in the following statements:}
		
		\begin{enumerate}[a)]
			\item Planets move in circular orbits around the sun, which is at the center.
			\vspace{35 pt}\\
			\item The sun, the center of our planetary system, does not move.
			\vspace{35 pt}\\
			\item A planet travels the same distances in its orbit around the sun at equal time intervals.
			\vspace{35 pt}\\
			\item If the Earth is at the aphelion of its orbit, then it is at its shortest distance from the Sun.
			\vspace{35 pt}\\
			\item American astronomer Edwin Hubble discovered a blue shift in the spectrum of stars and concluded that all galaxies are the same distance from Earth.
			\vspace{35 pt}\\
			\item If a cyclist rides at $\SI{18}{\kilo\meter\over \hour} $, then he travels $\SI{1.8}{\kilo\meter}$ in 10 minutes.
			\vspace{35 pt}\\
			\item At the center of the Milky Way, which stands still in the universe, is our sun.
		\end{enumerate}
		
	\end{enumerate}
	
	\fbox{
		\begin{minipage}{0.5\textwidth}
			For the solution please \href{https://www.okuyakl.de/phys/p10astL123/le123.pdf}{click here} or scan the QR code.\\
			Additional worksheets are available at
			
			\href{http://www.okuyakl.com}{www.okuyakl.com}
		\end{minipage}
		\hfill
		\begin{minipage}{0.4\textwidth}
			\includegraphics[width=1.5 cm]{../../viecher/zwe03}
			\includegraphics[width=3 cm]{qre123}
			\includegraphics[width=2 cm]{../../viecher/afanticon1}
			
	\end{minipage}}
	
\end{document}